\section{Problem Formulation and Foundations}
\label{sec:background}

\subsection{Planning Domains and Generalized Evidence}

We use typed STRIPS PDDL plus a bounded-integer resource fragment
\cite{McDermott1998PDDL}. A domain is
$D=\langle\mathcal{P},\mathcal{F},\mathcal{A},\mathcal{T}\rangle$, where
$\mathcal{P}$ and $\mathcal{F}$ contain predicate schemas and integer-valued
resource functions, $\mathcal{A}$ contains action schemas, and $\mathcal{T}$ is
a type hierarchy. A state $s=(s_{\mathcal P},s_{\mathcal F})$ combines Boolean
and integer valuations. For an applicable grounded action $a$, the Boolean
successor is
\[
\gamma_{\mathcal P}(s,a)=
(s_{\mathcal P}\setminus\mathit{del}(a))\cup\mathit{add}(a).
\]
The numeric component admits finite integer values, comparisons with constants,
and constant integer increases or decreases; arbitrary arithmetic and
continuous change are excluded. A generalized-planning task supplies training
and test instances that share $D$. MOOSE solves singleton goal conditions and
regresses the resulting plans into first-order condition-to-macro rules
\cite{Chen2026Moose}. We treat these rules as \emph{evidence}: they demonstrate
ways to achieve observed singleton goals, but need not constitute a closed or
minimal BDI plan library.

\subsection{BDI Plan Libraries and Their AgentSpeak(L) Realization}

A procedural BDI plan rule has a triggering event, an applicability context,
and a body. In the \asl{} notation used by our implementation
\cite{Rao1996AgentSpeak}, an achievement plan is
\[
+!p(\bar X):C\leftarrow b_1;\ldots;b_m,
\]
where $C$ is a conjunction of beliefs and comparisons and each $b_i$ is a PDDL
action or an internal achievement call $!q(\bar Y)$. A library is \emph{lifted}
when plans contain variables rather than training objects, and \emph{closed}
when every internal call has a selected implementation. Our certificates apply
to this trigger--context--body model and to the PDDL semantics of primitive
actions. We render the selected rules in \asl{} and execute them with Jason
\cite{Bordini2007Jason}; another BDI language would require its own renderer and
execution study. Query controllers also follow declarative goal patterns by
rechecking their intended state after subsidiary plans return
\cite{Hubner2006DeclarativeGoals}.

\subsection{Finite-Trace Temporal Goals}

An \ltlf{} formula is interpreted over a finite state sequence
\cite{DeGiacomoVardi2013}. We construct its deterministic finite automaton (DFA)
with \texttt{ltlf2dfa} and MONA \cite{Fuggitti2020LTLf2DFA}. The implemented
fragment admits certified distance-reducing transitions whose guards contain
conjunction and literal negation. A guard
\[
B=P^+\land\neg P^-
\]
contains achievements $P^+$ and atoms $P^-$ that must remain absent. A
query-local monitor evaluates the DFA initially and after every successful
primitive PDDL action. Disjunction, arbitrary arithmetic, and progress without
a certified action strategy are outside the implemented fragment.

\subsection{Certified Compilation Problem}

\textbf{Compilation problem.} Given $D$, training instances
$\mathcal{I}_{\mathrm{train}}$, singleton-goal evidence $E$, and a lifted query
$\varphi$, compilation produces one domain library $\mathcal{L}_D$ of lifted
atomic modules and query-local plans $\mathcal{Q}_{\varphi}$ that pursue
certified DFA progress. Positive predicate goals and supported integer
equalities require an executable achieving module. Negative literals denote
absence and preservation, never synthetic negative achievement calls.
Conjunctive guards require a preservation-safe order; mixed Boolean--numeric
guards require complete net effects for a common final state. The compiler
rejects unsupported arithmetic, disjunction, cyclic threats without a
preserving strategy, and any obligation whose effects are incomplete. In the
current rendering, only PDDL fluents and actions, static type facts
\texttt{obj\_tp/2}, and query-local controller names may appear.
